\section{Local return without a stationary density floor}
\label{sec:floor-free}
Let $\mu$ be uniform on the $L$-cube, $P=P_{1/2}$, $Q=P_{1/\sqrt2}$,
$h>0$, $\mathbb E hPh=1$, $g=Ph$, $\pi=hg$, and $U=h/g$.
The Doob transition is $K(x,y)=P(x,y)h(y)/g(x)$.
Suppose $h$ factors through a binary linear quotient with a fractional
basis packing of fixed weight $b\ge1$. The identity quotient with $b=1$
allows every positive potential. Write
\[
 a_b=1+2^b,\qquad r_b=a_b/(a_b-2).
\]
Assume $\|\pi\|_p\le B$ for fixed finite $p>r_b$, and that every
nonconstant stationary Walsh moment through order $k$ vanishes, $1\le k\le L$.

\paragraph{Theorem.}
For every sufficiently small fixed $\sigma>0$ depending on $b,p$,
put $A=\{h/(Ph)>e^{\sigma k}\}$. Let $s$ be the top eigenvalue of the
chain killed outside $A$, and $R_A$ the largest stationary one-step
return probability to a nonempty subset of $A$, conditional on starting
in that subset. Both are zero for empty $A$. For
$\Delta=\max_j|\lambda_j(K)-\lambda_j(P)|$ in decreasing order,
\[
 s/4-Ce^{-ck}\le\Delta\le s+Ce^{-ck},\qquad
 R_A\le s\le\sqrt{2R_A-R_A^2}.                                      \tag{FF1}
\]
The constants depend only on $B,b,p,\sigma$. The sorted union of the
ordinary spectrum and the killed spectrum, truncated to its first $2^L$
entries, recovers the full ordered list with multiplicities to this
accuracy. Thus $\Delta\to0$ precisely when $R_A\to0$ as $k\to\infty$.
This applies to general half-noise for $p>3$ and to two disjoint bases
for $p>5/3$, including $p=15/7$. The stationary floor in Section
\ref{sec:local-return-test} is removed by changing the exceptional set.
No inverse moment is required.

\paragraph{Truncating the actual multiplier.}
On quotient functions $P_\lambda$ maps $L^u$ contractively into
$L^{1+(u-1)\lambda^{-2b}}$ by the allocation of noise to bases above.
For a nonnegative nonzero quotient weight $W$ the unique positive
Perron vector of $QM_WQ$ is invariant under kernel translations.
The quotient estimate and H\"older therefore yield
$\|QM_WQ\|\le\|W\|_{r_b}$ on the ambient cube. Absolute values handle
signed weights by quadratic-form domination.

Fix $\lambda$ slightly above $1/\sqrt2$ and below one. Since $p>r_b$,
we can ensure $1+\lambda^{-2b}\ge2p/(p-1)$ and
$1+(p-1)\lambda^{-2b}\ge2$. Put $\theta=1/(\sqrt2\lambda)<1$.
Write $F=Qh$, $a=\mathbb E h$ and $V=\|F-a\|_2^2=1-a^2$.
Cauchy--Schwarz in $h\,d\mu$ gives $1\le a\,\mathbb E(\pi g)$.
Since $P_\lambda(\pi-1)$ has degrees above $k$,
\[
 V\le a|\langle P_\lambda(\pi-1),P_\theta(F-a)\rangle|
       \le (B+1)\theta^{k+1}\sqrt V .
\]
Thus $\|g-1\|_2\le C_B\theta^k$. Choose $0<\eta\le-\log(\theta)/2$.
The set $D=\{|g-1|>e^{-\eta k}\}$ has $\mu(D)\le C_B e^{-2\eta k}$.
Put $M=e^{\sigma k}$ and $U_0=U{\bf1}_{A^c}$. On $D$ the $L^{r_b}$
norms of $U_0$ and $\pi$ are at most $M\mu(D)^{1/r_b}$ and
$B\mu(D)^{1/r_b-1/p}$, respectively. Choose $\sigma<2\eta/r_b$.
Off $D$, for large $k$, $|U-\pi|\le Ce^{-\eta k}\pi$, and $U>M$
implies $\pi>M/2$. The $L^{r_b}$ norm of $\pi{\bf1}_{\pi>M/2}$
is at most $B^{p/r_b}(M/2)^{1-p/r_b}$. Altogether
$\|U_0-\pi\|_{r_b}\le Ce^{-ck}$.
At $\lambda$ the same Perron argument bounds
$\|P_\lambda M_{\pi-1}P_\lambda\|\le B+1$.
Sandwich by $P_\theta$ and split both factors at Walsh degree
$\lfloor k/2\rfloor$. Matching kills the low--low form and each remaining
block gains $\theta^{\lfloor k/2\rfloor+1}$. Hence
\[
 \|QM_UQ-P-S\|\le Ce^{-ck},\qquad S=QM_{U{\bf1}_A}Q .                \tag{FF2}
\]

\paragraph{A fourth-moment test from Markov contraction.}
The two Gram matrices of $M_{\sqrt U}Q$ have the same spectrum and
$M_{\sqrt U}PM_{\sqrt U}$ is the reversible conjugate of $K$.
Consequently $0\le QM_UQ\le I$. Testing on the constant gives
$\mathbb E U\le1$. Since $\mathbb E\pi=1$ and $U_0$ is exponentially
close to $\pi$ in $L^1$, $\mathbb E[U{\bf1}_A]\le Ce^{-ck}$.
For an ambient polynomial $f$ of Walsh degree at most $d$, $\|f\|_2\le1$,
apply the same contraction to $Q^{-1}f^2$:
\[
 \mathbb E[U f^4]\le\|Q^{-1}f^2\|_2^2\le2^{2d}\|f\|_4^4\le36^d .
\]
The final inequality is ordinary Boolean hypercontractivity. Therefore
$\mathbb E[U{\bf1}_A f^2]\le6^d(\mathbb E[U{\bf1}_A])^{1/2}$,
independently of the multiplicity of the low modes.

Take $d=\lfloor\delta k\rfloor$ for sufficiently small fixed $\delta>0$,
and let $\Pi$ project through degree $d$. Use $f=Qz$ for unit $z$ in
this subspace to obtain $\|\Pi S\Pi\|\le Ce^{-c'k}$. Since $0\le S\le I$,
the off-diagonal block has norm at most $\|\Pi S\Pi\|^{1/2}$.
Write $S_H$ for the compression to the complementary subspace. There
$P$ has norm at most $2^{-(d+1)}$. By (FF2), $QM_UQ$ is exponentially
close to the direct sum of the low part of $P$ and $S_H$.
The nonzero spectrum of $S$ equals that of $K_{AA}$ by the restricted
Gram identity. On two copies of the space $P\oplus S$ differs to the
same accuracy from the preceding direct sum padded with zeros.
Min--max proves the ordered reconstruction.

For the lower bound set $s_H=\|S_H\|=s+O(e^{-c'k})$. If
$s_H>4\,2^{-d}$, choose $1\le j\le d$ with $s_H/4<2^{-j}\le s_H/2$.
At the first index of level $2^{-j}$, the high block adds one
eigenvalue at least $2^{1-j}$ to all ordinary levels above $2^{-j}$.
This raises that index by at least $2^{-j}$. Smaller $s_H$ is absorbed
in the error. The upper bound follows from (FF2). The indicator and
level-set argument of Section \ref{sec:local-return-test} gives the two
inequalities involving $R_A$ using only reversibility.
The remaining critical question is whether all subsets of this
high-multiplier set have vanishing return for arbitrary two-base quotients.
